% This document was generated by an LLM.

\documentclass{essay}

\usepackage{centernot}

\title{The First and Second Derivative Tests for Local Extrema}
\author{}
\date{}

\begin{document}

\maketitle
\tableofcontents

\section{Introduction}

One of the central problems in elementary calculus is to find and classify the local extrema of a real-valued function. Given a function \(f : D \to \mathbb{R}\), we want to identify points at which \(f\) has a local maximum or local minimum, and, when possible, determine the type of extremum from information about the derivatives of \(f\).

Two standard tools are the \emph{first derivative test} and the \emph{second derivative test}. They are often presented as computational recipes:

\begin{enumerate}
    \item find the critical points by solving \(f'(x)=0\);
    \item inspect \(f'\) or \(f''\);
    \item declare a maximum or minimum.
\end{enumerate}

This summary is useful computationally, but it conceals several logical subtleties. A critical point need not be an extremum. An extremum need not occur at a point where \(f'(c)=0\). The second derivative test can fail to classify a genuine extremum. A point where \(f''(c)>0\) is a local minimum only under additional hypotheses. Endpoints require separate treatment. Points where the derivative does not exist may still be extrema.

The purpose of this essay is to state the two tests carefully, explain why they work, identify their assumptions, and examine examples and counterexamples that show exactly what can and cannot be concluded from derivative information. Since the second derivative test is often explained geometrically, we will also introduce concavity and inflection points: these are the concepts behind the informal language of ``bending upward'' and ``bending downward.''

\section{Local extrema and critical points}

Let \(c\) belong to the domain of \(f\).

We say that \(f\) has a \emph{local maximum} at \(c\) if there exists \(\delta>0\) such that
\[
f(x)\leq f(c)
\]
for every \(x\in D\) satisfying \(|x-c|<\delta\).

Similarly, \(f\) has a \emph{local minimum} at \(c\) if there exists \(\delta>0\) such that
\[
f(x)\geq f(c)
\]
for every \(x\in D\) satisfying \(|x-c|<\delta\).

A local maximum or minimum is called a \emph{local extremum}.

A point \(c\) in the domain of \(f\) is commonly called a \emph{critical point} or \emph{critical number} if either \(f'(c)=0\) or \(f'(c)\) does not exist.

Terminology varies somewhat between textbooks, but the important conceptual point is that critical points are \emph{candidates} for extrema. They are not automatically extrema.

For example, \(f(x)=x^3\) satisfies \(f'(0)=0\), but \(x=0\) is neither a local maximum nor a local minimum.

On the other hand, \(f(x)=|x|\) has a local minimum at \(x=0\), although \(f'(0)\) does not exist.

Thus the equation \(f'(c)=0\) is neither a definition of a local extremum nor a complete characterization of local extrema.

\section{Fermat's theorem: a necessary condition}

The derivative tests are closely related to an important necessary condition.

\medskip

\noindent
\textbf{Fermat's theorem.}
Suppose that \(f\) has a local extremum at an interior point \(c\) of its domain and that \(f\) is differentiable at \(c\). Then
\[
f'(c)=0.
\]

\medskip

The requirement that \(c\) be an interior point is essential. For example, \(f(x)=x\) on \([0,1]\) has an absolute and local minimum relative to its domain at \(x=0\), but \(f'(0)=1\) if the derivative is understood by extending \(f(x)=x\) to a neighborhood of \(0\).

The differentiability assumption is also essential. The function \(f(x)=|x|\) has a local minimum at \(0\), but \(f'(0)\) does not exist.

Fermat's theorem therefore says:

\[
\text{differentiable interior local extremum}
\quad\Longrightarrow\quad
f'(c)=0.
\]

The converse is false:
\[
f'(c)=0
\quad\centernot\Longrightarrow\quad
\text{\(c\) is a local extremum}.
\]

The derivative tests provide additional hypotheses that sometimes allow us to turn a critical point into a classified extremum.

\section{The first derivative test}

\subsection{Statement}

Suppose that \(f\) is continuous at \(c\) and differentiable on an interval immediately to the left and immediately to the right of \(c\), except possibly at \(c\) itself.

Then the sign of \(f'\) determines whether \(f\) is increasing or decreasing on each side of \(c\).

The first derivative test can be stated as follows.

\begin{enumerate}
    \item If \(f'(x)>0\) for \(x<c\) sufficiently close to \(c\), and \(f'(x)<0\) for \(x>c\) sufficiently close to \(c\), then \(f\) has a local maximum at \(c\).

    \item If \(f'(x)<0\) for \(x<c\) sufficiently close to \(c\), and \(f'(x)>0\) for \(x>c\) sufficiently close to \(c\), then \(f\) has a local minimum at \(c\).

    \item If \(f'\) has the same sign on both sides of \(c\), then \(f\) has no local extremum at \(c\).
\end{enumerate}

In compact form,
\[
+\longrightarrow -
\quad\Rightarrow\quad
\text{local maximum},
\]
whereas
\[
-\longrightarrow +
\quad\Rightarrow\quad
\text{local minimum}.
\]

The crucial information is not merely that \(f'(c)=0\). The crucial information is a \emph{change in the sign of \(f'\)} across \(c\).

\section{Why the first derivative test works}

Suppose, for example, that \(f'(x)<0\) to the left of \(c\), while \(f'(x)>0\) to the right of \(c\).

The sign of the derivative implies that \(f\) is decreasing as \(x\) approaches \(c\) from the left and increasing as \(x\) moves away from \(c\) to the right. Consequently, nearby values of \(f\) lie above \(f(c)\).

More formally, choose \(x<c\) sufficiently close to \(c\). By the Mean Value Theorem, assuming the usual continuity and differentiability hypotheses on \([x,c]\), there exists \(\xi\in(x,c)\) such that
\[
\frac{f(c)-f(x)}{c-x}=f'(\xi).
\]
Since \(f'(\xi)<0\) and \(c-x>0\), we obtain \(f(c)-f(x)<0\), so \(f(c)<f(x)\).

Likewise, for \(x>c\), the Mean Value Theorem gives a point \(\eta\in(c,x)\) such that
\[
\frac{f(x)-f(c)}{x-c}=f'(\eta)>0.
\]
Hence \(f(x)>f(c)\).

Thus \(f(c)<f(x)\) for all sufficiently nearby \(x\neq c\), which means that \(f\) has a strict local minimum at \(c\).

The proof for a local maximum is analogous.

This exposes an important hidden ingredient: the usual justification for the test relies on the theorem that the sign of the derivative controls monotonicity, which in turn normally follows from the Mean Value Theorem.

\section{Examples of the first derivative test}

\subsection{A local minimum}

Consider \(f(x)=x^2\). Then \(f'(x)=2x\). For \(x<0\), \(f'(x)<0\), while for \(x>0\), \(f'(x)>0\).
Thus
\[
-\longrightarrow +
\]
at \(x=0\), so \(f\) has a local minimum there.

\subsection{A local maximum}

Consider \(f(x)=-x^2\). Then \(f'(x)=-2x\).
For \(x<0\), \(f'(x)>0\), while for \(x>0\), \(f'(x)<0\). Hence
\[
+\longrightarrow -
\]
and \(x=0\) is a local maximum.

\subsection{A stationary point that is not an extremum}

Consider \(f(x)=x^3\). Then \(f'(x)=3x^2\). Although \(f'(0)=0\), we have \(f'(x)>0\) on both sides of \(0\). Therefore the function is increasing through \(0\), and \(0\) is neither a local maximum nor a local minimum.

This is one of the most important counterexamples to the incorrect implication
\[
f'(c)=0
\quad\Longrightarrow\quad
\text{\(c\) is an extremum}.
\]

\subsection{An extremum where the derivative does not exist}

Consider \(f(x)=|x|\). For \(x<0\), \(f'(x)=-1\), while for \(x>0\), \(f'(x)=1\).
Thus the sign changes from negative to positive, and \(x=0\) is a local minimum.

Notice that \(f'(0)\) does not exist. The first derivative test can therefore classify certain extrema at nondifferentiable points. This is one advantage it has over the usual second derivative test.

\section{Failure modes of the first derivative test}

The first derivative test is logically strong when its hypotheses can be verified, but it may be unusable or inconclusive in several situations.

\subsection{No sign change}

If \(f'\) has the same sign on both sides of a critical point, the point is not an extremum.

For \(f(x)=x^3\), we have \(f'(x)=3x^2>0\) for every \(x\neq0\). The zero derivative at the origin does not cause an extremum.

\subsection{The derivative changes sign infinitely often}

Consider a function such as
\[
f(x)=
\begin{cases}
x^4\sin(1/x), & x\neq0,\\
0, & x=0.
\end{cases}
\]
Near \(0\), its derivative has oscillatory behavior. There may be no single interval on each side of \(0\) on which the derivative keeps a fixed sign.

The simple sign-chart version of the first derivative test is then not directly applicable.

\subsection{Failure of continuity at the candidate point}

It is dangerous to infer an extremum merely from derivative signs on the punctured neighborhood if the value \(f(c)\) is unrelated to the nearby values.

For example, define
\[
f(x)=
\begin{cases}
x^2, & x\neq0,\\
1, & x=0.
\end{cases}
\]
For \(x<0\), \(f'(x)=2x<0\), and for \(x>0\), \(f'(x)=2x>0\).
The derivative sign pattern resembles that of a local minimum.

But \(f(0)=1\), while nearby values \(x^2\) are much smaller than \(1\). Therefore \(0\) is actually a strict local maximum, not a local minimum.

The point is that derivative information on the two punctured sides controls the behavior of the function on those sides, but it does not by itself determine the isolated value \(f(c)\). Continuity at \(c\) connects \(f(c)\) with the surrounding behavior.

\subsection{Endpoints}

At an endpoint, there may be only one side of the domain. The ordinary two-sided first derivative test does not apply.

For example, \(f(x)=x\) on \([0,1]\) has a minimum at \(0\) and a maximum at \(1\), although there is no sign change of \(f'\) across either endpoint.

When optimizing on a closed interval, endpoints must therefore be examined separately.

\section{Concavity}

Before turning to the second derivative test, it is useful to separate two different kinds of shape information.

The first derivative describes whether the function itself is increasing or decreasing.

The second derivative describes whether the first derivative is increasing or decreasing. Geometrically, this is information about the way the graph bends.

\subsection{Geometric intuition from tangent lines}

The most useful geometric picture is this: concave up means the graph lies above its tangent lines.

At a point \(a\), the tangent line has equation
\[
y=f(a)+f'(a)(x-a).
\]
So, for a differentiable function, concave up behavior means that the tangent line at \(a\) is a supporting line underneath the graph:
\[
f(x)\geq f(a)+f'(a)(x-a).
\]

This captures the idea of bending upward more rigorously than the phrase ``cup-shaped.'' The graph may rise or fall, and the tangent line need not be horizontal. What matters is that, after touching the graph at \(a\), the tangent line stays below the graph nearby, or on the whole interval when the function is concave up on an interval.

For example, \(f(x)=x^2\) is concave up. Its derivative \(f'(x)=2x\) increases as \(x\) increases.

Similarly, concave down means the graph lies below its tangent lines.

For a differentiable function, concave down behavior means
\[
f(x)\leq f(a)+f'(a)(x-a).
\]

The tangent line is now a supporting line above the graph.

For example, \(f(x)=-x^2\) is concave down. Its derivative \(f'(x)=-2x\) decreases as \(x\) increases.

This is why concavity is naturally connected with the second derivative: \(f''\) is the derivative of \(f'\), so it measures the rate at which the slope is changing.

\subsection{Definitions using chords}

Let \(I\) be an interval, and let \(f:I\to\mathbb R\).

For functions that may not be differentiable, concavity is usually defined using chords rather than tangent lines.

The function \(f\) is \emph{concave up} on \(I\), also called \emph{convex} on \(I\), if for every \(x,y\in I\) and every \(t\in[0,1]\),
\[
f(tx+(1-t)y)
\leq
t f(x)+(1-t)f(y).
\]

Geometrically, this says that the graph of \(f\) lies below or on each chord joining two points of the graph.

The function \(f\) is \emph{concave down} on \(I\) if for every \(x,y\in I\) and every \(t\in[0,1]\),
\[
f(tx+(1-t)y)
\geq
t f(x)+(1-t)f(y).
\]

Geometrically, this says that the graph of \(f\) lies above or on each chord joining two points of the graph.

The inequalities are non-strict. This matters. A line is both concave up and concave down, because its graph agrees exactly with every chord.

\subsection{Definitions using tangent lines}

When \(f\) is differentiable on an interval, the chord definition is equivalent to the tangent-line inequalities above.

If \(f\) is concave up on an interval, then its graph lies above or on its tangent lines:
\[
f(x)\geq f(a)+f'(a)(x-a)
\]
for \(a,x\) in the interval.

If \(f\) is concave down on an interval, then its graph lies below or on its tangent lines:
\[
f(x)\leq f(a)+f'(a)(x-a)
\]
for \(a,x\) in the interval.

This tangent-line interpretation is often the most useful one for extrema. If the tangent line at \(c\) is horizontal, then the tangent line has equation \(y=f(c)\). If the graph is concave up and lies above this horizontal tangent line, then nearby values of \(f\) satisfy \(f(x)\geq f(c)\), which is exactly the shape of a local minimum. If the graph is concave down and lies below the horizontal tangent line, then nearby values satisfy \(f(x)\leq f(c)\), which is exactly the shape of a local maximum.

\subsection{Derivative criteria for concavity}

The connection with derivatives is as follows.

If \(f\) is differentiable on an interval \(I\), then \(f\) being concave up on \(I\) corresponds to \(f'\) being nondecreasing on \(I\).

Similarly, \(f\) being concave down on \(I\) corresponds to \(f'\) being nonincreasing on \(I\).

If \(f\) is twice differentiable on \(I\), then the usual second-derivative criteria are:
\[
f''(x)\geq0\text{ on }I
\quad\Longrightarrow\quad
f\text{ is concave up on }I,
\]
and
\[
f''(x)\leq0\text{ on }I
\quad\Longrightarrow\quad
f\text{ is concave down on }I.
\]

In elementary calculus, this is often shortened informally to:
\[
f''>0 \text{ means concave up},
\qquad
f''<0 \text{ means concave down}.
\]

The interval language is important. Concavity is a property of behavior over an interval, not just a property of a single point. A statement such as \(f''(c)>0\) is pointwise information; by itself it is not the same thing as saying that \(f\) is concave up on a whole neighborhood of \(c\).

\subsection{Inflection points}

An \emph{inflection point} is a point at which the graph changes concavity.

More explicitly, \(c\) is an inflection point if, on one side of \(c\), the function is concave up, while on the other side it is concave down, or conversely.

For example, \(f(x)=x^3\) has \(f''(x)=6x\). For \(x<0\), \(f''(x)<0\), so the graph is concave down. For \(x>0\), \(f''(x)>0\), so the graph is concave up. Thus \(0\) is an inflection point.

Notice that \(0\) is not a local extremum of \(x^3\). An inflection point is about a change in bending, not about being highest or lowest nearby.

This distinction is central to the derivative tests: local extrema concern the sign pattern of \(f'\), whereas inflection points concern the concavity pattern of \(f\).

\section{The second derivative test}

\subsection{Standard statement}

Suppose that \(f\) is twice differentiable near an interior point \(c\), and suppose that \(f'(c)=0\).

Then:

\begin{enumerate}
    \item if \(f''(c)>0\), then \(f\) has a strict local minimum at \(c\);

    \item if \(f''(c)<0\), then \(f\) has a strict local maximum at \(c\);

    \item if \(f''(c)=0\), the test is inconclusive.
\end{enumerate}

The last statement is especially important. It does \emph{not} say that \(c\) is not an extremum. It says that the second derivative test alone gives no classification.

\section{Why the second derivative test works}

The second derivative measures the local rate of change of the first derivative.

Suppose \(f'(c)=0\) and \(f''(c)>0\).
By the definition of the second derivative,
\[
f''(c)
=
\lim_{x\to c}
\frac{f'(x)-f'(c)}{x-c}.
\]
Since \(f'(c)=0\),
\[
f''(c)
=
\lim_{x\to c}
\frac{f'(x)}{x-c}.
\]

If this limit is positive, then for \(x\) sufficiently close to \(c\),
\[
\frac{f'(x)}{x-c}>0.
\]

When \(x<c\), the denominator \(x-c\) is negative, so \(f'(x)<0\).

When \(x>c\), the denominator \(x-c\) is positive, so \(f'(x)>0\).

Therefore
\[
-\longrightarrow +
\]
for the sign of \(f'\), and the first derivative test gives a local minimum.

Similarly, if \(f''(c)<0\), then the sign pattern becomes
\[
+\longrightarrow -,
\]
and \(c\) is a local maximum.

This shows that the second derivative test is not an entirely separate phenomenon. It is, in a precise sense, a convenient sufficient condition that forces the sign change required by the first derivative test.

\section{Concavity interpretation of the second derivative test}

The geometric interpretation can now be stated more precisely.

If \(f\) is concave up near \(c\), then its graph lies above its tangent lines. If, in addition, \(f'(c)=0\), then the tangent line at \(c\) is horizontal: \(y=f(c)\). The tangent-line picture therefore suggests \(f(x)\geq f(c)\) near \(c\), so \(c\) should be a local minimum.

Similarly, if \(f\) is concave down near \(c\) and \(f'(c)=0\), then the graph lies below the horizontal tangent line, suggesting \(f(x)\leq f(c)\) near \(c\), so \(c\) should be a local maximum.

This is the geometric version of the second derivative test.

There are two cautions.

First, concavity is interval-level information. Saying that the graph is concave up on a whole interval around \(c\) is stronger than knowing only the single number \(f''(c)\). In many elementary examples, \(f''\) is continuous, so \(f''(c)>0\) implies \(f''(x)>0\) near \(c\), and therefore the graph is concave up near \(c\). But the standard proof of the second derivative test does not need to pass through an interval concavity theorem; it can use the limit definition of \(f''(c)\) directly.

Second, positive second derivative by itself does not imply a minimum unless the first derivative vanishes at the point.

For example, \(f(x)=x^2+x\) has \(f''(0)=2>0\), but \(x=0\) is not a local minimum because \(f'(0)=1\neq0\).

The actual minimum occurs at \(x=-\frac12\).

Thus the implication
\[
f''(c)>0
\quad\Longrightarrow\quad
\text{local minimum}
\]
is false.

The correct version requires, among other things, \(f'(c)=0\).

\section{Taylor expansion as a local model}

Taylor's formula gives another way to understand why the derivative tests have the form they do.

Suppose \(f\) is sufficiently smooth near an interior point \(c\). For \(h\) close to \(0\), Taylor's formula says that
\[
f(c+h)
=
f(c)
+
f'(c)h
+
\frac12 f''(c)h^2
+
\text{higher-order terms}.
\]
Equivalently,
\[
f(c+h)-f(c)
=
f'(c)h
+
\frac12 f''(c)h^2
+
\text{higher-order terms}.
\]

This expansion is a local model for the change in \(f\) near \(c\).

\subsection{The linear term}

If \(f'(c)\neq0\), then the linear term \(f'(c)h\) dominates for sufficiently small \(h\). But \(h\) is negative on the left of \(c\) and positive on the right of \(c\). Therefore the expression \(f'(c)h\) has opposite signs on the two sides of \(c\).

This is the Taylor-expansion reason behind Fermat's theorem. If \(f'(c)\neq0\), then nearby values of \(f\) lie both above and below \(f(c)\), so \(c\) cannot be an interior local maximum or local minimum.

\subsection{The quadratic term}

If \(f'(c)=0\), then the linear term disappears and the expansion begins as
\[
f(c+h)-f(c)
=
\frac12 f''(c)h^2
+
\text{higher-order terms}.
\]

If \(f''(c)>0\), then the leading quadratic term is positive for every sufficiently small \(h\neq0\), because \(h^2>0\).
This models a strict local minimum.

If \(f''(c)<0\), then the leading quadratic term is negative for every sufficiently small \(h\neq0\). This models a strict local maximum.

This is the Taylor-expansion picture behind the second derivative test. At a stationary point, the sign of the first nonzero quadratic term determines whether nearby values initially sit above or below \(f(c)\).

\subsection{Connection with the first derivative test}

Taylor expansion also explains why the second derivative test forces the sign pattern used by the first derivative test.

Apply the same local-model idea to \(f'\). If \(f'\) is differentiable near \(c\), then
\[
f'(c+h)
=
f'(c)
+
f''(c)h
+
\text{higher-order terms}.
\]

If \(f'(c)=0\) and \(f''(c)>0\), then \(f'(c+h)\) has the sign of \(h\) for small nonzero \(h\). Thus \(f'\) is negative to the left of \(c\) and positive to the right of \(c\), which is the first derivative test pattern
\[
-\longrightarrow +
\]
for a local minimum.

If \(f'(c)=0\) and \(f''(c)<0\), then \(f'(c+h)\) has the opposite sign from \(h\). Thus \(f'\) is positive to the left of \(c\) and negative to the right of \(c\), which is the pattern
\[
+\longrightarrow -
\]
for a local maximum.

\subsection{The role of the remainder}

Taylor expansion is an explanation, but it is not magic. To use it as a proof, one needs hypotheses that control the remainder term.

For example, if
\[
f(c+h)-f(c)
=
\frac12 f''(c)h^2
+
o(h^2)
\]
and \(f''(c)>0\), then the \(o(h^2)\) term is eventually smaller in magnitude than the positive quadratic term. This makes the sign conclusion rigorous.

Without a remainder estimate, the phrase ``higher-order terms'' is only a heuristic placeholder. This is why the earlier proof of the second derivative test used the definition of \(f''(c)\) directly: it gave exactly the sign control needed near \(c\).

\section{Examples of the second derivative test}

\subsection{Positive second derivative}

Let \(f(x)=x^2+3\).
Then
\[
f'(x)=2x,
\qquad
f''(x)=2.
\]
At \(c=0\), \(f'(0)=0\) and \(f''(0)=2>0\).
Therefore \(0\) is a strict local minimum.

\subsection{Negative second derivative}

Let \(f(x)=5-x^2\).
Then
\[
f'(x)=-2x,
\qquad
f''(x)=-2.
\]
At \(0\), \(f'(0)=0\) and \(f''(0)=-2<0\).
Therefore \(0\) is a strict local maximum.

\section{The most important failure mode: \texorpdfstring{\(f''(c)=0\)}{f''(c)=0}}

Suppose \(f'(c)=0\) and \(f''(c)=0\).
Nothing definite follows.

All three basic possibilities can occur.

\subsection{A local minimum}

Consider \(f(x)=x^4\).
Then
\[
f'(x)=4x^3,
\qquad
f''(x)=12x^2.
\]
At \(0\), \(f'(0)=0\) and \(f''(0)=0\).
Nevertheless,
\[
x^4\geq0=x^4\big|_{x=0},
\]
so \(0\) is a strict local minimum.

\subsection{A local maximum}

Consider \(f(x)=-x^4\). Again, \(f'(0)=0\) and \(f''(0)=0\), but \(0\) is a strict local maximum.

\subsection{Neither maximum nor minimum}

Consider \(f(x)=x^3\). We have \(f'(0)=0\) and \(f''(0)=0\), but \(0\) is neither a local maximum nor a local minimum.

In fact, \(0\) is an inflection point: the graph changes from concave down for \(x<0\) to concave up for \(x>0\). This reinforces the distinction between a change in concavity and a local extremum.

Therefore the data \(f'(c)=0\) and \(f''(c)=0\) are compatible with all three outcomes:
\[
\text{minimum},\qquad
\text{maximum},\qquad
\text{neither}.
\]

This is exactly what it means for the second derivative test to be \emph{inconclusive}.

\section{A stronger counterexample: an extremum with no useful second derivative}

The function \(f(x)=|x|\) has a strict local minimum at \(0\), but neither \(f'(0)\) nor \(f''(0)\) exists.

Thus the second derivative test is not a necessary criterion for extrema. It is only a sufficient criterion under its hypotheses.

A less sharp example is \(f(x)=|x|^{3/2}\). This function is differentiable at \(0\), with \(f'(0)=0\), and \(0\) is a strict local minimum, but the second derivative is singular near \(0\). Again, classification by the standard second derivative test is unavailable.

\section{The distinction between necessary and sufficient conditions}

Much confusion about derivative tests disappears once the relevant statements are classified logically.

For an interior point \(c\), assuming differentiability at \(c\),
\[
\text{\(c\) is a local extremum}
\quad\Longrightarrow\quad
f'(c)=0.
\]
This is a \emph{necessary condition} supplied by Fermat's theorem.

But \(f'(c)=0\) is not sufficient.

Under suitable neighborhood hypotheses,
\[
f'(x)<0\text{ to the left of }c
\quad\text{and}\quad
f'(x)>0\text{ to the right of }c
\]
is sufficient for a local minimum.

Likewise, the conditions \(f'(c)=0\) and \(f''(c)>0\) are sufficient for a strict local minimum under the hypotheses of the second derivative test.

But neither the first nor the second derivative test gives a necessary characterization of every possible local extremum.

The logical structure is therefore asymmetric:
\[
\text{extremum}
\Longrightarrow
\text{critical candidate}
\]
under suitable differentiability and interior-point assumptions, while additional derivative behavior may imply
\[
\text{critical candidate}
\Longrightarrow
\text{extremum}.
\]

\section{Strict and non-strict extrema}

The derivative tests naturally detect \emph{strict} extrema when the derivative actually changes sign.

However, local extrema need not be strict.

Consider the constant function \(f(x)=0\).
Every point is simultaneously a local maximum and a local minimum according to the usual non-strict definitions, since
\(f(x)=f(c)\) for every \(x\).

Yet \(f'(x)=0\) everywhere, so there is no sign change in the derivative.

Another example is a function with a flat interval, such as
\[
f(x)=
\begin{cases}
x^2, & x<0,\\
0, & 0\leq x\leq1,\\
(x-1)^2, & x>1.
\end{cases}
\]
Every point of \([0,1]\) is a local minimum, but an interior point of the flat interval does not exhibit the usual negative-to-positive derivative sign change.

Therefore the first derivative test should not be misunderstood as an ``if and only if'' characterization of all non-strict extrema.

\section{Endpoints and constrained domains}

Suppose \(f:[a,b]\to\mathbb{R}\).
An absolute maximum or minimum can occur at

\begin{enumerate}
    \item an interior critical point;
    \item an interior point where the derivative does not exist;
    \item the endpoint \(a\);
    \item the endpoint \(b\).
\end{enumerate}

For this reason, the standard closed-interval method is not merely ``solve \(f'(x)=0\).'' Instead one evaluates the function at all relevant critical points and at both endpoints, then compares the resulting values.

For example, \(f(x)=x\) on \([0,1]\) has no interior critical point, but it has \(\min_{[0,1]}f=f(0)=0\) and \(\max_{[0,1]}f=f(1)=1\).

The derivative tests concern local behavior around interior points. Optimization on a restricted domain always requires attention to the geometry of the domain itself.

\section{Higher derivative tests}

When the second derivative test is inconclusive, higher derivatives sometimes resolve the situation.

Suppose \(f'(c)=f''(c)=\cdots=f^{(n-1)}(c)=0\) and \(f^{(n)}(c)\neq0\).

Under appropriate smoothness assumptions:

\begin{enumerate}
    \item if \(n\) is even and \(f^{(n)}(c)>0\), then \(c\) is a local minimum;

    \item if \(n\) is even and \(f^{(n)}(c)<0\), then \(c\) is a local maximum;

    \item if \(n\) is odd, then \(c\) is not a local extremum.
\end{enumerate}

This extends the Taylor-expansion viewpoint from the second derivative test. The leading nonzero term controls the local sign of \(f(x)-f(c)\). In Taylor form,
\[
f(x)
=
f(c)
+
\frac{f^{(n)}(c)}{n!}(x-c)^n
+
\text{higher-order terms}.
\]

If \(n\) is even, the sign of \((x-c)^n\) is the same on both sides of \(c\). If \(n\) is odd, it changes sign.

For example, \(f(x)=x^4\) has its first nonzero derivative at order \(4\), which is even and positive, so \(0\) is a local minimum.

For \(f(x)=x^3\), the first nonzero derivative occurs at order \(3\), which is odd, so \(0\) is not an extremum.

\section{But even all finite derivative tests may fail}

Higher derivatives do not provide a universal solution.

Consider
\[
f(x)=
\begin{cases}
e^{-1/x^2}, & x\neq0,\\
0, & x=0.
\end{cases}
\]
This function is infinitely differentiable, and \(f^{(n)}(0)=0\) for every positive integer \(n\).

Nevertheless, \(f(x)>0=f(0)\) for every \(x\neq0\), so \(0\) is a strict local minimum.

Thus all derivatives at the point can vanish while the point is still a strict local minimum.

This example shows that even complete knowledge of the finite-order derivatives \emph{at the point itself} need not determine the local geometry of a smooth function.

The function is not represented near \(0\) by its Taylor series. Its Taylor series at \(0\) is simply \(0\), even though the function is positive away from \(0\).

This illustrates a deeper distinction between smooth functions and analytic functions.

\section{The first derivative test versus the second derivative test}

The two tests are related but have different strengths.

The first derivative test examines the behavior of \(f'\) on a neighborhood of the candidate point. It asks whether \(f\) changes from increasing to decreasing or from decreasing to increasing.

The second derivative test uses more localized information. If \(f'(c)=0\) and \(f''(c)\neq0\), the sign of \(f''(c)\) already forces the required change in \(f'\).

Thus the second derivative test is often computationally simpler, but it applies in fewer situations.

For example, \(f(x)=x^4\) is immediately classified by the first derivative test: \(f'(x)=4x^3\) changes from negative to positive at \(0\).

The second derivative test fails because \(f''(0)=0\).

Similarly, \(f(x)=|x|\) can be classified by a one-sided derivative sign analysis, but the ordinary second derivative test cannot even be formulated at \(0\).

In this sense, the first derivative test is often more robust.

\section{A useful hierarchy of questions}

When analyzing extrema of a one-variable function, it is useful to proceed in the following logical order.

\begin{enumerate}
    \item \textbf{What is the domain?}

    Identify endpoints, excluded points, and other domain restrictions before differentiating.

    \item \textbf{Where are the candidate points?}

    Look for points where \(f'(x)=0\) and points in the domain where \(f'\) does not exist. If the domain is restricted, include relevant endpoints.

    \item \textbf{Can the first derivative test be applied?}

    Determine the sign of \(f'\) on the two sides of each interior candidate.

    \item \textbf{Can the second derivative test be applied more efficiently?}

    At a stationary point \(c\), compute \(f''(c)\). If it is nonzero, classification follows immediately.

    \item \textbf{If \(f''(c)=0\), do not conclude anything yet.}

    Return to the first derivative test, examine the function directly, use higher derivatives when justified, or use another local argument.

    \item \textbf{For absolute extrema on a closed interval, compare values.}

    Local classification alone does not determine which local extremum is the absolute one.
\end{enumerate}

\section{Common incorrect inferences}

Several errors occur frequently.

\subsection{\texorpdfstring{``\(f'(c)=0\), therefore \(c\) is an extremum''}{``f'(c)=0, therefore c is an extremum''}}

False. The counterexample is \(f(x)=x^3\).

\subsection{\texorpdfstring{``\(f''(c)>0\), therefore \(c\) is a minimum''}{``f''(c)>0, therefore c is a minimum''}}

False unless the required stationary-point hypothesis is also present.

For example, \(f(x)=x^2+x\) satisfies \(f''(0)>0\), but \(0\) is not a local minimum.

\subsection{\texorpdfstring{``\(f''(c)=0\), therefore \(c\) is an inflection point''}{``f''(c)=0, therefore c is an inflection point''}}

False.

For \(f(x)=x^4\), we have \(f''(0)=0\), but \(0\) is a local minimum, not an inflection point in the usual concavity-change sense.

The more relevant test is whether concavity changes across the point. For \(g(x)=x^3\), we have \(g''(x)=6x\), which changes sign at \(0\). That sign change detects a change from concave down to concave up. For \(x^4\), by contrast, \(f''(x)=12x^2\geq0\) on both sides of \(0\), so there is no change in concavity.

\subsection{``If the second derivative test fails, there is no extremum''}

False.

Again, \(f(x)=x^4\) has a strict local minimum at \(0\), despite \(f''(0)=0\).

\subsection{\texorpdfstring{``Every extremum occurs where \(f'(c)=0\)''}{``Every extremum occurs where f'(c)=0''}}

False.

The function \(f(x)=|x|\) has a minimum at \(0\), where the derivative does not exist.

Endpoints provide another family of counterexamples.

\subsection{``A critical point is an extremum''}

False. A critical point is merely a candidate requiring further analysis.

\section{A compact comparison}

The main facts may be summarized as follows.

\[
\begin{array}{c|c|c}
\text{Information} & \text{Conclusion} & \text{Status}\\
\hline
f'(c)=0
& \text{no definite conclusion}
& \text{candidate only}\\[4pt]
f':+\to-
& \text{local maximum}
& \text{first derivative test}\\[4pt]
f':-\to+
& \text{local minimum}
& \text{first derivative test}\\[4pt]
f' \text{ has same sign on both sides}
& \text{no local extremum}
& \text{first derivative test}\\[4pt]
f''\geq0 \text{ on an interval}
& \text{concave up}
& \text{concavity criterion}\\[4pt]
f''\leq0 \text{ on an interval}
& \text{concave down}
& \text{concavity criterion}\\[4pt]
f'' \text{ changes sign across } c
& \text{inflection point candidate}
& \text{concavity change}\\[4pt]
f'(c)=0,\ f''(c)>0
& \text{strict local minimum}
& \text{second derivative test}\\[4pt]
f'(c)=0,\ f''(c)<0
& \text{strict local maximum}
& \text{second derivative test}\\[4pt]
f'(c)=0,\ f''(c)=0
& \text{no definite conclusion}
& \text{test inconclusive}
\end{array}
\]

This table should not be used without remembering the regularity and domain assumptions behind the statements.

\section{Conclusion}

The derivative tests are not definitions of extrema. They are theorems that connect local derivative behavior with local order behavior of the function.

The first derivative test is fundamentally about monotonicity. A local minimum occurs when the function changes from decreasing to increasing, and a local maximum occurs when it changes from increasing to decreasing. The derivative provides a convenient way to detect these changes.

Concavity is about a different question: how the graph bends, or equivalently how the slopes change. Concave up means that the slopes are nondecreasing; concave down means that the slopes are nonincreasing. Inflection points are places where this concavity changes.

The second derivative test is a more specialized shortcut connecting these ideas. At a stationary point, a nonzero second derivative forces the appropriate sign change of the first derivative: \(f''(c)>0\) forces a negative-to-positive transition and hence a local minimum, while \(f''(c)<0\) forces a positive-to-negative transition and hence a local maximum.

Geometrically, the same conclusion says: if the tangent line is horizontal and the graph bends upward in the relevant local sense, the point is a minimum; if the tangent line is horizontal and the graph bends downward, the point is a maximum.

Taylor expansion gives the same idea algebraically: the first nonzero term in the local expansion of \(f(c+h)-f(c)\) describes the first visible shape of the graph near \(c\).

But the logic is deliberately one-sided. The conditions \(f'(c)=0\) and \(f''(c)>0\) are sufficient for a strict local minimum, not necessary. A minimum may occur when \(f''(c)=0\), when \(f''(c)\) does not exist, or even when \(f'(c)\) does not exist. Likewise, the equation \(f'(c)=0\) identifies only a candidate, not an extremum.

The most important habit is therefore to distinguish three questions:

\begin{enumerate}
    \item Why is this point a candidate?
    \item Which theorem is being used to classify it?
    \item Are the hypotheses of that theorem actually satisfied?
    \item Is the question about an extremum, or about a change in concavity?
\end{enumerate}

Once these questions are kept separate, the first and second derivative tests stop being mechanical rules and become examples of a broader pattern in analysis: local information about derivatives can imply qualitative information about a function, but only under precisely stated assumptions.

\end{document}
